
@@M0@@
\uline{On
suppose que} @@M1@@ \uline{est quasi-compact et a un faisceau ample}. \uline{Alors}, \uline{on a la suite exacte}:
@@M2@@
\uline{où on pose}, \uline{pour} @@M3@@:
@@M4@@
\uline{et pour} @@M5@@, @@M6@@:
@@M7@@
\uline{les notations étant celles de} 3.1 \uline{et} 3.3. \uline{En outre}, @@M8@@ \uline{admet un inverse à gauche} @@M9@@, \uline{défini par} @@M10@@.

\uline{Si de plus} @@M11@@ \uline{est noethérien}, \uline{et} @@M12@@ \uline{régulier}, \uline{alors} @@M13@@ \uline{est surjectif}, \uline{donc on a la suite exacte}:
@@M14@@

Tout d'abord, @@M15@@ admet l'inverse à gauche @@M16@@, et a fortiori est injectif. En effet, on a:
@@M17@@
d'après la formule de projection. Par ailleurs, @@M18@@ (VI 5.9), d'où le résultat.

Ensuite, @@M19@@. Si @@M20@@, on a en effet:
@@M21@@
d'après 3.4.

Il reste donc à montrer que @@M22@@. Soient donc @@M23@@, @@M24@@, tels que:
